%%%%%%%%%%%%%%%%%%%%%%%%%%%%%%%%%%%%%%%%
%
% $ Autor: Gruppe 07 $
% $ Projekt: Magic Faucet Fountain $
% $ Pfad: LaTeXTemplate/Nano33BLESense/Contents/General/Hardwarebeschreibung/DokuProjektMagicFaucetFountain/BeschreibungSystem.tex $
% $ Kurzbeschreibung: Dieses Dokument stellt das Projekt Magic Faucet Fountain vor. Es erklärt den Aufbau, die Verdrahtung und Programmierung der Applikation. $
%
% !TeX spellcheck = de_DE/GB
% !TeX program = pdflatex
% !BIB program = biber/bibtex
% !TeX encoding = utf8
%
%%%%%%%%%%%%%%%%%%%%%%%%%%%%%%%%%%%%%%%%


\chapter{Projekt: Magic Faucet Fountain}


\section{Beschreibung}
Das Projekt Magic Faucen Fountain beschäftigt sich mit der Umsetzung eines Brunnens, der die optische Täuschung eines schwebenden Wasserhahns erzeugt. Dabei scheint es, als würde Wasser kontinuierlich aus einem frei schwebenden Hahn fließen - ganz ohne sichtbare Verbindung zur Wasserquelle \ref{MagicFaucet}.
\begin{figure}[H]
	\centering
	\caption{Skizze des Magic Faucet Fountain}
	\resizebox{0.5\textwidth}{!}{%
		\begin{circuitikz}
			\tikzstyle{every node}=[font=\LARGE]
			\draw [line width=1.2pt, short] (3.5,13.25) -- (6,3.25);
			\draw [line width=1.2pt, short] (22.5,13.25) -- (20,3.25);
			\draw [line width=1.2pt, short] (6,3.25) .. controls (9.75,0.75) and (16.25,0.75) .. (20,3.25);
			\draw [line width=1.2pt, short] (12.5,12.25) -- (12.5,26.75);
			\draw [line width=1.2pt, short] (13.5,12.25) -- (13.5,26.75);
			\draw [line width=1.2pt, short] (6.25,12) .. controls (6,13.25) and (10.75,13.5) .. (12.5,13.5);
			\draw [line width=1.2pt, short] (13.5,13.5) .. controls (17.75,13.5) and (20.25,12.75) .. (19.75,12);
			\draw [line width=1.2pt, short] (12.5,15.25) .. controls (0.5,14.75) and (2.25,11.25) .. (13,11.25);
			\draw [line width=1.2pt, short] (13.5,15.25) .. controls (24.25,15) and (24.75,11.25) .. (13,11.25);
			\draw [line width=1.2pt, short] (12.5,15.75) .. controls (2.25,15.5) and (-1.25,11.25) .. (13,10.75);
			\draw [line width=1.2pt, short] (13.5,15.75) .. controls (26.5,15) and (24.5,10.75) .. (13,10.75);
			\draw [line width=1.2pt, short] (12.5,12.25) .. controls (12.75,12) and (13.25,12) .. (13.5,12.25);
			\draw [ line width=1.2pt ] (19.5,7.25) circle (0.25cm);
			\draw [line width=1.2pt, short] (12,27) -- (12,27.5);
			\draw [line width=1.2pt, short] (14,27) -- (14,27.5);
			\draw [line width=1.2pt, short] (12.5,27.5) .. controls (12.25,30) and (13,30.5) .. (14.5,31.25);
			\draw [line width=1.2pt, short] (13.5,27.5) .. controls (13.25,29) and (13.5,29.5) .. (14.5,30);
			\draw [line width=1.2pt, short] (14.5,31.25) .. controls (15,31.25) and (15.25,31.5) .. (15.25,32);
			\draw [line width=1.2pt, short] (17.25,32) -- (17.25,32.5);
			\draw [line width=1.2pt, short] (16,32.5) -- (16,33.5);
			\draw [line width=1.2pt, short] (16.5,32.5) -- (16.5,33.5);
			\draw [ line width=1.2pt ] (16.25,34.75) ellipse (0.5cm and 0.25cm);
			\draw [line width=1.2pt, short] (15.75,34.75) -- (14.25,34.5);
			\draw [line width=1.2pt, short] (14.5,34) -- (16,34.25);
			\draw [line width=1.2pt, short] (15.75,34.25) -- (15.75,33.75);
			\draw [line width=1.2pt, short] (16.75,34.75) -- (16.75,33.75);
			\draw [line width=1.2pt, short] (16.75,34.5) -- (18,34.75);
			\draw [line width=1.2pt, short] (16.5,35) -- (17.75,35.25);
			\draw [line width=1.2pt, short] (14.25,34.5) .. controls (14.5,34.25) and (14.75,34.25) .. (14.5,34);
			\draw [line width=1.2pt, short] (14.25,34.5) .. controls (14,34) and (14.25,34) .. (14.5,34);
			\draw [line width=1.2pt, short] (17.75,35.25) .. controls (18,35) and (18.25,35) .. (18,34.75);
			\draw [line width=1.2pt, short] (15.75,33.75) .. controls (16,33.25) and (16.75,33.5) .. (16.75,33.75);
			\draw [line width=1.2pt, short] (16,32.75) .. controls (15,32.5) and (15,32.25) .. (16.25,32.25);
			\draw [line width=1.2pt, short] (16.5,32.75) .. controls (17.5,32.5) and (17.5,32.25) .. (16.25,32.25);
			\draw [line width=1.2pt, short] (15.25,32.5) -- (15.25,32);
			\draw [line width=1.2pt, short] (15.25,32) .. controls (15.75,31.5) and (17,31.75) .. (17.25,32);
			\draw [line width=1.2pt, short] (16,32.5) .. controls (16,32.25) and (16.5,32.25) .. (16.5,32.5);
			\draw [line width=1.2pt, short] (15,29.75) .. controls (16,29.25) and (16.25,29) .. (17.5,29.5);
			\draw [line width=1.2pt, short] (17.25,32) .. controls (18.25,31.5) and (18.75,30) .. (17.5,29.5);
			\draw [line width=1.2pt, short] (14.5,30) .. controls (14.5,30) and (14.75,30) .. (15,29.75);
			\draw [line width=1.2pt, short] (17.75,31.75) -- (19,32.25);
			\draw [line width=1.2pt, short] (18.25,30.75) -- (19.5,31.25);
			\draw [line width=1.2pt, short] (19,32.25) .. controls (19.25,32.25) and (19.75,31.25) .. (19.5,31.25);
			\draw [line width=1.2pt, short] (19,32.75) -- (19.75,33);
			\draw [line width=1.2pt, short] (19.75,30.75) -- (20.5,31);
			\draw [line width=1.2pt, short] (19,32.75) .. controls (20,32.75) and (20.5,30.75) .. (19.75,30.75);
			\draw [line width=1.2pt, short] (19,31) .. controls (19.25,30.75) and (19.5,30.5) .. (19.75,30.75);
			\draw [line width=1.2pt, short] (19,32.75) .. controls (19,32.5) and (18.75,32.5) .. (18.75,32.25);
			\draw [line width=1.2pt, short] (19.75,33) .. controls (20.5,32.75) and (21,31.5) .. (20.5,31);
			\draw [line width=1.2pt, short] (12.5,27.75) .. controls (11.25,27.25) and (12.5,27) .. (13,27.25);
			\draw [line width=1.2pt, short] (13.5,27.75) .. controls (14.75,27.25) and (13.5,27) .. (13,27.25);
			\draw [line width=1.2pt, short] (12,27) .. controls (12.25,26.75) and (12.5,26.75) .. (13,26.75);
			\draw [line width=1.2pt, short] (13,26.75) .. controls (13.5,26.75) and (13.75,26.75) .. (14,27);
			\draw [line width=1.2pt, short] (12.5,27.5) .. controls (12.75,27.25) and (13.25,27.25) .. (13.5,27.5);
		\end{circuitikz}
	}%
	\label{MagicFaucet}
\end{figure}

Dieser Effekt wird durch ein transparentes Acrylrohr realisiert, das sowohl den Wasserfluss leitet als auch die tragende Struktur für den Hahn darstellt. Da das Rohr hinter dem gleichmäßig fließenden Wasser optisch kaum wahrnehmbar ist, entsteht der Eindruck eines schwebenden Hahns.

Der Wasserstrom wird über eine kleine elektrische Pumpe erzeugt, die über ein Relais angesteuert wird. Das Relais fungiert als elektronischer Schalter und wird von einem digitalen Ausgang des Arduino-Boards gesteuert. Der Arduino ist über Bluetooth mit einer mobilen App verbunden. Über die App kann der Benutzer den Brunnen bequem ein- und ausschalten.

Ein weiterer Bestandteil des Systems ist ein Ultraschallsensor, der den Wasserstand im Vorratsbehälter misst. Der aktuelle Füllstand wird ebenfalls in der App angezeigt, wodurch der Benutzer rechtzeitig erkennen kann, wann Wasser nachgefüllt werden muss.


\section{Aufbau}
Die mechanische Struktur des Magic Faucet Fountain basiert auf einem dekorativen Blumentopf, der gleichzeitig als Wasserreservoir und Trägerelement für den sichtbaren Aufbau dient. Der Blumentopf wird durch einen funktional gestalteten Deckel verschlossen, der sowohl mechanische als auch gestalterische Funktionen erfüllt.

Zentral im Deckel befindet sich eine Durchgangsbohrung, durch die ein transparentes Acrylrohr senkrecht nach oben geführt wird. Am oberen Ende des Rohrs ist ein dekorativer Wasserhahn angebracht. Direkt unterhalb dieses Hahns befinden sich feine Austrittsbohrungen, durch die Wasser kontrolliert herausfließt. Aufgrund der Transparenz des Rohrs und des gleichmäßigen Wasserstroms entsteht die Illusion eines frei schwebenden, endlos fließenden Wasserhahns.

Seitlich im Deckel befinden sich zusätzliche Rücklaufbohrungen, über die das Wasser zurück in den Blumentopf gelangt. Eine im Inneren platzierte Wasserpumpe sorgt für den geschlossenen Wasserkreislauf, indem sie das Wasser wieder nach oben in das Acrylrohr befördert.

Die elektronischen und steuerungstechnischen Komponenten wie Arduino, Relais, Spannungsversorgung sowie die Status-LEDs und der Kippschalter befinden sich in einem separaten Gehäuse, das als elektronischer Kasten (E-Kasten) ausgeführt ist. Dieser Kasten schützt die Elektronik vor Feuchtigkeit und erlaubt einen strukturierten, servicefreundlichen Aufbau. Die Leitungen zur Pumpe, zum Ultraschallsensor und zu den LEDs werden durch entsprechende Kabeldurchführungen zwischen dem E-Kasten und der Hauptbaugruppe verbunden.

Zur optischen Aufwertung kann die Oberseite des Deckels mit dekorativen Steinen oder anderen Gestaltungselementen versehen werden, wodurch die technische Funktion harmonisch in ein ästhetisches Gesamtbild eingebettet wird.

\subsection{Abmaße}
Da der endgültige Aufbau des Projekts „Magic Faucet Fountain“ zum Zeitpunkt der Dokumentation noch nicht erfolgt ist, können die Abmessungen lediglich auf Basis der entworfenen Komponenten und Gehäusekonzepte geschätzt werden. Die Steuerungseinheit inklusive Arduino, Relais-Schaltung und Stromversorgung soll in einem kompakten Gehäuse untergebracht werden. Die Wasserinstallation mit Pumpe und transparentem Rohrsystem wird so dimensioniert, dass sie sich auf einer Fläche von ca. 40 cm × 40 cm bei einer Höhe von maximal 80 cm aufbauen lässt.

\section{Hardware-Schnittstellen}
Die Magic Faucet Fountain besteht aus verschiedenen Hardware-Komponenten, welche durch physische Verbindungen miteinander kommunizieren und Daten austauschen. In diesem Kapitel sollen diese Schnittstellen mit Schaltplänen veranschaulicht mittels dieser erklärt werden.

\subsection{Schaltpläne}
\subsubsection{Versorgungsspannungen 5 V und 12 V}
Die im Schaltplan in Abbildung \ref{SchaltplanVersorgung} dargestellten Komponenten umfassen den Netzanschluss (230 V AC / 50 Hz) sowie das Schaltnetzteil vom Typ MW RD-50A. Dieses wandelt die Netzspannung in geregelte Gleichspannungen von 5 V und 12 V DC um, die für den Betrieb der Steuereinheit und der Aktoren erforderlich sind. Der Anschluss erfolgt über:
\begin{itemize}
	\item 3-adriges Netzkabel (\(1{,}5\,\text{mm}^2\))
	\item 4 Verbindungskabel (\(0{,}5\,\text{mm}^2\)) für:
	\begin{itemize}
		\item +5 V (z. B. Arduino, BT-Modul)
		\item GND (5 V)
		\item +12 V (z. B. Pumpe)
		\item GND (12 V)
	\end{itemize}
\end{itemize}
\begin{figure}[H]
	\centering
	\caption{Versorgungsschaltplan}
	\label{SchaltplanVersorgung}
	\resizebox{1\textwidth}{!}{%
		\begin{circuitikz}
			\tikzstyle{every node}=[font=\Huge]
			\draw [ line width=2pt](-27.5,-1) to[short, -o] (-27.5,-2.25) ;
			\draw [ line width=2pt](-26.25,-1) to[short, -o] (-26.25,-2.25) ;
			\draw [ line width=2pt](-25,-1) to[short, -o] (-25,-2.25) ;
			\draw [ line width=2pt , rounded corners = 12.0] (-12.5,6.5) rectangle (-2.5,1.5);
			\draw [ line width=2pt](-11.25,5.5) to[short, -o] (-11.25,6.75) ;
			\draw [ line width=2pt](-10,5.5) to[short, -o] (-10,6.75) ;
			\draw [ line width=2pt](-8.75,5.5) to[short, -o] (-8.75,6.75) ;
			\node [font=\Huge, rotate around={90:(0,0)}] at (-27.5,-3.75) {\textbf{L1}};
			\node [font=\Huge, rotate around={90:(0,0)}] at (-26.25,-3.5) {\textbf{N}};
			\node [font=\Huge, rotate around={90:(0,0)}] at (-25,-3.5) {\textbf{PE}};
			\draw [ line width=2pt , rounded corners = 12.0, dashed] (-28.75,-1) rectangle  (-23.75,-4.75);
			\node [font=\Huge] at (-26.5,-5.75) {\textbf{Netzkabel}};
			\node [font=\Huge] at (-26.5,-7) {\textbf{230 V / 50 Hz}};
			\node [font=\Huge] at (-11.25,4.5) {\textbf{1}};
			\node [font=\Huge] at (-10,4.5) {\textbf{2}};
			\node [font=\Huge] at (-8.75,4.5) {\textbf{3}};
			\draw [ color={rgb,255:red,255; green,0; blue,0}, line width=2pt, short] (-27.5,-0.75) -- (-27.5,12.75);
			\draw [ color={rgb,255:red,255; green,0; blue,0}, line width=2pt, short] (-27.5,12.75) -- (-11.25,12.75);
			\draw [ color={rgb,255:red,255; green,0; blue,0}, line width=2pt, short] (-11.25,12.75) -- (-11.25,6.75);
			\draw [ color={rgb,255:red,0; green,0; blue,255}, line width=2pt, short] (-26.25,-0.75) -- (-26.25,11.5);
			\draw [ color={rgb,255:red,0; green,0; blue,255}, line width=2pt, short] (-26.25,11.5) -- (-10,11.5);
			\draw [ color={rgb,255:red,0; green,0; blue,255}, line width=2pt, short] (-10,11.5) -- (-10,6.75);
			\draw [ color={rgb,255:red,0; green,128; blue,0}, line width=2pt, short] (-25,-0.75) -- (-25,10.25);
			\draw [ color={rgb,255:red,0; green,128; blue,0}, line width=2pt, short] (-25,10.25) -- (-8.75,10.25);
			\draw [ color={rgb,255:red,0; green,128; blue,0}, line width=2pt, short] (-8.75,10.25) -- (-8.75,7);
			\draw [ line width=2pt](-7.5,5.5) to[short, -o] (-7.5,6.75) ;
			\draw [ line width=2pt](-6.25,5.5) to[short, -o] (-6.25,6.75) ;
			\draw [ line width=2pt](-5,5.5) to[short, -o] (-5,6.75) ;
			\draw [ line width=2pt](-3.75,5.5) to[short, -o] (-3.75,6.75) ;
			\node [font=\Huge, color={rgb,255:red,0; green,0; blue,255}] at (-8.25,2.5) {\textbf{MW RD-50A}};
			\draw [ line width=2pt](-7.5,7) to[short] (-7.5,12.75);
			\draw [ line width=2pt](-7.5,12.75) to[short] (3.75,12.75);
			\draw [ color={rgb,255:red,255; green,128; blue,0}, , line width=2pt](-6.25,6.75) to[short] (-6.25,11.5);
			\draw [ color={rgb,255:red,255; green,128; blue,0}, , line width=2pt](-6.25,11.5) to[short] (3.75,11.5);
			\draw [ line width=2pt](-5,7) to[short] (-5,9);
			\draw [ line width=2pt](-5,9) to[short] (3.75,9);
			\draw [ color={rgb,255:red,128; green,0; blue,255}, , line width=2pt](-3.75,6.75) to[short] (-3.75,7.75);
			\draw [ color={rgb,255:red,128; green,0; blue,255}, , line width=2pt](-3.75,7.75) to[short] (3.75,7.75);
			\node [font=\Huge, color={rgb,255:red,128; green,0; blue,255}] at (5.25,7.75) {\textbf{+ 12 V}};
			\node [font=\Huge] at (6.25,9) {\textbf{GND (12V)}};
			\node [font=\Huge] at (6,12.75) {\textbf{GND (5V)}};
			\node [font=\Huge, color={rgb,255:red,255; green,128; blue,0}] at (5.25,11.5) {\textbf{+ 5 V}};
			\node [font=\Huge] at (-7.5,4.5) {\textbf{4}};
			\node [font=\Huge] at (-6.25,4.5) {\textbf{5}};
			\node [font=\Huge] at (-5,4.5) {\textbf{6}};
			\node [font=\Huge] at (-3.75,4.5) {\textbf{7}};
			\draw [ line width=2pt ] (-26.25,4) ellipse (2.5cm and 0.25cm);
			\node [font=\Huge] at (-20.75,4) {\textbf{3 x $\emptyset$1,5 mm2}};
			\draw [line width=2pt, dashed] (1.25,7.25) -- (1.25,13.25);
			\node [font=\Huge] at (1.25,13.75) {\textbf{4 x $\emptyset$0,5 mm2}};
		\end{circuitikz}
	}%
\end{figure}

\subsubsection{5 V Netz mit Arduino und Ultraschallsensor}
Im Schaltplan \ref{5VSchaltplan} sind der Arduino Nano, der Ultraschallsensor und die Stromversorgung durch das Netzteil als Komponenten zu entnehmen. Es lassen sich folgende Schnittstellen aus der Schaltung herleiten:
	\begin{itemize}
	\item Digitale Ausgänge D2, D3, D4 zur Steuerung der Sig-Anschlüsse der Relais für die Pumpe und LEDs
	\item Spannungsversorgung 5 V / GND
\end{itemize}

\begin{figure}[H]
	\centering
	\caption{5V - Netz}
	\label{5VSchaltplan}
	\resizebox{1\textwidth}{!}{%
\begin{circuitikz}
	\tikzstyle{every node}=[font=\Huge]
	\draw [ color={rgb,255:red,255; green,128; blue,0}, line width=2pt, short] (-8.5,9) -- (0,9);
	\draw [ color={rgb,255:red,255; green,128; blue,0}, line width=2pt, short] (2.75,9) -- (10.25,9);
	\draw [ line width=2pt , rounded corners = 12.0, dashed] (-2.5,12.75) rectangle  (5,7.75);
	\draw [line width=2pt, short] (0,9) -- (2.5,10.25);
	\draw [line width=2pt, short] (1.25,9.75) -- (1.25,11.5);
	\draw [line width=2pt, short] (0,11.5) -- (2.5,11.5);
	\draw [line width=2pt, short] (0,11.5) -- (0,11);
	\draw [line width=2pt, short] (2.5,11.5) -- (2.5,11);
	\node [font=\LARGE] at (-0.25,10) {3};
	\node [font=\LARGE] at (3.25,10) {4};
	\node [font=\Huge, color={rgb,255:red,0; green,0; blue,255}] at (1.25,13.5) {\textbf{Kippschalter}};
	\node [font=\Huge, color={rgb,255:red,255; green,128; blue,0}] at (-10.25,9) {\textbf{+ 5 V}};
	\draw [ color={rgb,255:red,255; green,128; blue,0}, line width=2pt, short] (10.25,9) -- (20,9);
	\draw [ color={rgb,255:red,255; green,128; blue,0}, line width=2pt, short] (20,9) -- (20,7.75);
	\draw [ color={rgb,255:red,255; green,128; blue,0}, line width=2pt, short] (20,7.75) -- (20,7.75);
	\draw [ color={rgb,255:red,255; green,128; blue,0}, line width=2pt, short] (20,7.75) -- (20,6.5);
	\draw [ color={rgb,255:red,255; green,128; blue,0}, line width=2pt, ->, >=Stealth] (18.75,9) -- (25,9);
	\draw [ color={rgb,255:red,255; green,128; blue,0}, line width=2pt, ->, >=Stealth] (20,7.75) -- (25,7.75);
	\node [font=\Huge] at (27.5,7.75) {\textbf{Relais.2/+}};
	\node [font=\Huge] at (27.5,6.5) {\textbf{Relais.3/+}};
	\node [font=\Huge] at (27.5,9) {\textbf{Relais.1/+
	}};
	\draw [ line width=2pt , rounded corners = 12.0] (6.25,1.5) rectangle (16.25,-2.25);
	\draw [ color={rgb,255:red,255; green,128; blue,0}, line width=2pt, ->, >=Stealth] (20,6.5) -- (25,6.5);
	\draw [ color={rgb,255:red,255; green,128; blue,0}, line width=2pt, short] (17.5,9) -- (17.5,-7.25);
	\draw [ line width=2pt](7.5,0.25) to[short, -o] (7.5,2) ;
	\draw [ line width=2pt](10,0.25) to[short, -o] (10,2) ;
	\draw [ line width=2pt](12.5,0.25) to[short, -o] (12.5,2) ;
	\draw [ line width=2pt](15,0.25) to[short, -o] (15,2) ;
	\draw [ color={rgb,255:red,255; green,128; blue,0}, , line width=2pt](7.5,2.25) to[short] (7.5,9);
	\node at (7.5,9) [circ, color={rgb,255:red,255; green,128; blue,0}] {};
	\node at (17.5,9) [circ, color={rgb,255:red,255; green,128; blue,0}] {};
	\node at (20,9) [circ, color={rgb,255:red,255; green,128; blue,0}] {};
	\node at (20,7.75) [circ, color={rgb,255:red,255; green,128; blue,0}] {};
	\draw [ line width=2pt , rounded corners = 12.0] (-3.75,-8.5) rectangle (16.25,-13.5);
	\draw [line width=2pt, short] (-8.75,-18.5) -- (22.5,-18.5);
	\draw [ line width=2pt](15,-12.75) to[short, -o] (15,-14.75) ;
	\draw [ line width=2pt](12.5,-12.75) to[short, -o] (12.5,-14.75) ;
	\draw [ line width=2pt](12.5,-18.5) to[short] (12.5,-15);
	\node at (12.5,-18.5) [circ] {};
	\draw [ color={rgb,255:red,255; green,128; blue,0}, , line width=2pt](17.5,-7.25) to[short] (17.5,-16);
	\draw [ color={rgb,255:red,255; green,128; blue,0}, , line width=2pt](17.5,-16) to[short] (15,-16);
	\draw [ color={rgb,255:red,255; green,128; blue,0}, , line width=2pt](15,-16) to[short] (15,-14.75);
	\node [font=\Huge] at (15,-12) {\textbf{Vin}};
	\node [font=\Huge] at (12.5,-12) {\textbf{GND}};
	\node [font=\Huge] at (-11.25,-18.5) {\textbf{GND (5V)}};
	\draw [ line width=2pt](18.75,3.25) to[short] (18.75,-16.75);
	\draw [ line width=2pt](18.75,-16.5) to[short] (18.75,-18.5);
	\node at (18.75,-18.5) [circ] {};
	\node [font=\Huge] at (7.5,-0.25) {Vcc};
	\node [font=\Huge] at (9.75,-0.25) {Trig};
	\node [font=\Huge] at (12.5,-0.25) {Echo};
	\node [font=\Huge] at (15,-0.25) {GND};
	\draw [ line width=2pt](18.75,3.25) to[short] (18.75,4);
	\draw [ line width=2pt](18.75,4) to[short] (15,4);
	\draw [ line width=2pt](15,4) to[short] (15,2.25);
	\draw [ line width=2pt](15,-9) to[short, -o] (15,-7.25) ;
	\draw [ line width=2pt](12.5,-9) to[short, -o] (12.5,-7.25) ;
	\draw [ line width=2pt](10,-9) to[short, -o] (10,-7.25) ;
	\draw [ line width=2pt](7.5,-9) to[short, -o] (7.5,-7.25) ;
	\draw [ line width=2pt](5,-9) to[short, -o] (5,-7.25) ;
	\node [font=\Huge] at (15,-9.75) {\textbf{D2}};
	\node [font=\Huge] at (12.5,-9.75) {\textbf{D3}};
	\node [font=\Huge] at (10,-9.75) {\textbf{D4}};
	\node [font=\Huge] at (7.5,-9.75) {\textbf{D6}};
	\node [font=\Huge] at (5,-9.75) {\textbf{D9}};
	\node [font=\Huge, color={rgb,255:red,0; green,0; blue,255}] at (2.75,-12) {\textbf{Arduino Nano 33 BLE Sense}};
	\draw [ color={rgb,255:red,0; green,255; blue,0}, , line width=2pt](10,2) to[short] (10,4);
	\draw [ color={rgb,255:red,0; green,255; blue,0}, , line width=2pt](10,4) to[short] (5,4);
	\draw [ color={rgb,255:red,0; green,255; blue,0}, , line width=2pt](5,4) to[short] (5,-7);
	\draw [ color={rgb,255:red,0; green,255; blue,0}, , line width=2pt](12.5,2) to[short] (12.5,5.25);
	\draw [ color={rgb,255:red,0; green,255; blue,0}, , line width=2pt](12.5,5.25) to[short] (6.25,5.25);
	\draw [ color={rgb,255:red,0; green,255; blue,0}, , line width=2pt](6.25,5.25) to[short] (3.75,5.25);
	\draw [ color={rgb,255:red,0; green,255; blue,0}, , line width=2pt](3.75,5.25) to[short] (3.75,-4.75);
	\draw [ color={rgb,255:red,0; green,255; blue,0}, , line width=2pt](3.75,-4.75) to[short] (7.5,-4.75);
	\draw [ color={rgb,255:red,0; green,255; blue,0}, , line width=2pt](7.5,-4.75) to[short] (7.5,-7);
	\draw [ color={rgb,255:red,128; green,128; blue,0}, , line width=2pt](15,-7.25) to[short] (15,-6);
	\draw [ color={rgb,255:red,128; green,128; blue,0}, , line width=2pt](15,-6) to[short] (25,-6);
	\draw [ color={rgb,255:red,128; green,128; blue,0}, , line width=2pt](12.5,-7.25) to[short] (12.5,-4.75);
	\draw [ color={rgb,255:red,128; green,128; blue,0}, , line width=2pt](12.5,-4.75) to[short] (25,-4.75);
	\draw [ color={rgb,255:red,128; green,128; blue,0}, , line width=2pt](10,-7) to[short] (10,-3.5);
	\draw [ color={rgb,255:red,128; green,128; blue,0}, , line width=2pt](10,-3.5) to[short] (25,-3.5);
	\draw [ color={rgb,255:red,128; green,128; blue,0}, line width=2pt, ->, >=Stealth] (23.75,-3.5) -- (25,-3.5);
	\draw [ color={rgb,255:red,128; green,128; blue,0}, line width=2pt, ->, >=Stealth] (23.75,-4.75) -- (25,-4.75);
	\draw [ color={rgb,255:red,128; green,128; blue,0}, line width=2pt, ->, >=Stealth] (23.75,-6) -- (25,-6);
	\node [font=\Huge] at (28.5,-3.5) {\textbf{Relais.1/LED Grün}};
	\node [font=\Huge] at (29,-4.75) {\textbf{Relais.2/LED.Blau}};
	\node [font=\Huge] at (29.25,-6) {\textbf{Relais.3/Pumpe}};
	\node [font=\Huge, color={rgb,255:red,0; green,0; blue,255}] at (11.25,-1.5) {\textbf{HC-SR04}};
	\draw [line width=2pt, short] (22.5,-18.5) -- (22.5,-16);
	\draw [line width=2pt, ->, >=Stealth] (22.5,-16) -- (25,-16);
	\draw [line width=2pt, ->, >=Stealth] (22.5,-14.75) -- (25,-14.75);
	\draw [line width=2pt, ->, >=Stealth] (22.5,-13.5) -- (25,-13.5);
	\draw [line width=2pt, short] (22.5,-16) -- (22.5,-13.5);
	\node at (22.5,-14.75) [circ] {};
	\node at (22.5,-16) [circ] {};
	\node [font=\Huge] at (27.5,-13.5) {\textbf{Relais.1/-}};
	\node [font=\Huge] at (27.5,-14.75) {\textbf{Relais.2/-}};
	\node [font=\Huge] at (27.5,-16) {\textbf{Relais.3/-}};
	\end{circuitikz}
	}%
\end{figure}

Der Arduino verarbeitet die über das Bluetooth-Modul eingehenden Steuerbefehle und schaltet gezielt die Ausgänge zur Ansteuerung der Aktoren. Über das integrierte BLE-Modul erfolgt die drahtlose Kommunikation mit dem Smartphone.



\subsubsection{5/12 V Netz mit Aktoren}
Komponenten dieser Schaltung sind: 3x 5V-Relais, LED-Leuchten (grün und blau) und die Wasserpumpe. Dem Aufbau der Schaltung ist zu entnehmen, dass jedes Relais als Schalter arbeitet, das über den Arduino gesteuert wird. Die genaue Ansteuerung der Komponenten durch den Arduino wird in Tabelle \ref{AnsteuerungArduino} dargestellt.
	\begin{center}
	\captionof{table}{Ansteuerung der Komponenten durch den Arduino}
	\begin{tabular}{c|c|c|c}
		MOS-FET & Gesteuertes Element & Arduino Pin & Versorgungsspannung\\ \hline
		 1 & LED grün & D4 & 5 V\\
	     2 & LED blau & D3 & 5 V\\
	     3 & Wasserpumpe & D2 & 12 V\\
	\end{tabular}
	\label{AnsteuerungArduino}
\end{center}
Eine detaillierte Darstellung der Schaltung ist in Abbildung \ref{12VSchaltplan} zu finden.
\begin{figure}[H]
	\centering
	\caption{5/12 V Netz mit Aktoren}
	\resizebox{1\textwidth}{!}{%
		\begin{circuitikz}
			\tikzstyle{every node}=[font=\Huge]
			\draw [ line width=2pt , rounded corners = 12.0] (1.25,4) rectangle (8.75,-6);
			\draw [ line width=2pt , rounded corners = 12.0] (16.25,4) rectangle (23.75,-6);
			\draw [ line width=2pt , rounded corners = 12.0] (31.25,4) rectangle (38.75,-6);
			\draw [ line width=2pt ] (5,-14.5) circle (3.5cm);
			\draw [ line width=2pt ] (35,-14.5) circle (3.5cm);
			\draw [ line width=2pt ] (20,-14.5) circle (3.5cm);
			\node [font=\Huge, color={rgb,255:red,0; green,0; blue,255}] at (5,-14.75) {\textbf{Pumpe}};
			\node [font=\Huge, color={rgb,255:red,0; green,0; blue,255}] at (20,-14.75) {\textbf{LED.Blau}};
			\node [font=\Huge, color={rgb,255:red,0; green,0; blue,255}] at (35,-14.75) {\textbf{LED.Grün}};
			\node [font=\Huge, color={rgb,255:red,0; green,0; blue,255}] at (5,-1.25) {\textbf{Relais.1}};
			\node [font=\Huge, color={rgb,255:red,0; green,0; blue,255}] at (20,-1) {\textbf{Relais.2}};
			\node [font=\Huge, color={rgb,255:red,0; green,0; blue,255}] at (35,-1) {\textbf{Relais.3}};
			\draw [ line width=2pt](1.25,-1) to[short, -o] (0,-1) ;
			\draw [ line width=2pt](1.25,-3.5) to[short, -o] (0,-3.5) ;
			\draw [ line width=2pt](1.25,1.5) to[short, -o] (0,1.5) ;
			\node [font=\Huge] at (2,-3.5) {\textbf{+}};
			\node [font=\Huge] at (2,-1) {\textbf{-}};
			\node [font=\Huge] at (2.25,1.5) {\textbf{Sig}};
			\draw [ line width=2pt](16.25,-3.5) to[short, -o] (15,-3.5) ;
			\draw [ line width=2pt](16.25,-1) to[short, -o] (15,-1) ;
			\draw [ line width=2pt](16.25,1.5) to[short, -o] (15,1.5) ;
			\draw [ line width=2pt](31.25,-3.5) to[short, -o] (30,-3.5) ;
			\draw [ line width=2pt](31.25,-1) to[short, -o] (30,-1) ;
			\draw [ line width=2pt](31.25,1.5) to[short, -o] (30,1.5) ;
			\node [font=\Huge] at (17.25,-3.5) {\textbf{+}};
			\node [font=\Huge] at (17.25,-1) {\textbf{-}};
			\node [font=\Huge] at (17.5,1.5) {\textbf{Sig}};
			\node [font=\Huge] at (32,-3.5) {\textbf{+}};
			\node [font=\Huge] at (32,-1) {\textbf{-}};
			\node [font=\Huge] at (32.25,1.5) {\textbf{Sig}};
			\draw [ color={rgb,255:red,255; green,128; blue,0}, line width=2pt, ->, >=Stealth] (-8.75,10.25) -- (-6.25,10.25);
			\draw [ color={rgb,255:red,255; green,128; blue,0}, line width=2pt, short] (-0.25,-3.5) -- (-5,-3.5);
			\draw [ color={rgb,255:red,255; green,128; blue,0}, line width=2pt, short] (-5,-3.5) -- (-5,10.25);
			\draw [ color={rgb,255:red,255; green,128; blue,0}, line width=2pt, short] (15,-3.5) -- (10,-3.5);
			\draw [ color={rgb,255:red,255; green,128; blue,0}, line width=2pt, short] (10,-3.5) -- (10,11.5);
			\draw [ color={rgb,255:red,255; green,128; blue,0}, line width=2pt, short] (30,-3.5) -- (25,-3.5);
			\draw [ color={rgb,255:red,255; green,128; blue,0}, line width=2pt, short] (25,-3.5) -- (25,10.25);
			\draw [ color={rgb,255:red,255; green,128; blue,0}, line width=2pt, ->, >=Stealth] (-8.75,11.5) -- (-6.25,11.5);
			\draw [ color={rgb,255:red,255; green,128; blue,0}, line width=2pt, ->, >=Stealth] (-8.75,12.75) -- (-6.25,12.75);
			\draw [line width=2pt, ->, >=Stealth] (-8.75,-22.25) -- (-6.25,-22.25);
			\draw [line width=2pt, ->, >=Stealth] (-8.75,-23.5) -- (-6.25,-23.5);
			\draw [line width=2pt, ->, >=Stealth] (-8.75,-24.75) -- (-6.25,-24.75);
			\draw [line width=2pt, short] (-6.25,-22.25) -- (-5,-22.25);
			\draw [line width=2pt, short] (-5,-22.25) -- (-3.75,-22.25);
			\draw [line width=2pt, short] (-3.75,-22.25) -- (-3.75,-1);
			\draw [line width=2pt, short] (-3.75,-1) -- (-0.25,-1);
			\draw [line width=2pt, short] (-6.25,-23.5) -- (11.25,-23.5);
			\draw [line width=2pt, short] (11.25,-23.5) -- (11.25,-1);
			\draw [line width=2pt, short] (11.25,-1) -- (14.75,-1);
			\draw [line width=2pt, short] (-6.25,-24.75) -- (26.25,-24.75);
			\draw [line width=2pt, short] (26.25,-24.5) -- (26.25,-1);
			\draw [line width=2pt, short] (26.25,-1) -- (29.75,-1);
			\draw [ color={rgb,255:red,128; green,128; blue,0}, line width=2pt, ->, >=Stealth] (-8.75,7.75) -- (-6.25,7.75);
			\draw [ color={rgb,255:red,128; green,128; blue,0}, line width=2pt, ->, >=Stealth] (-8.75,6.5) -- (-6.25,6.5);
			\draw [ color={rgb,255:red,128; green,128; blue,0}, line width=2pt, ->, >=Stealth] (-8.75,5.25) -- (-6.25,5.25);
			\draw [ color={rgb,255:red,128; green,128; blue,0}, line width=2pt, short] (-2.5,1.5) -- (-0.25,1.5);
			\draw [ color={rgb,255:red,128; green,128; blue,0}, line width=2pt, short] (-6.25,7.75) -- (-2.5,7.75);
			\draw [ color={rgb,255:red,128; green,128; blue,0}, line width=2pt, short] (-2.5,7.75) -- (-2.5,1.5);
			\draw [ color={rgb,255:red,128; green,128; blue,0}, line width=2pt, short] (-6.25,6.5) -- (13.75,6.5);
			\draw [ color={rgb,255:red,128; green,128; blue,0}, line width=2pt, short] (13.75,6.5) -- (13.75,1.5);
			\draw [ color={rgb,255:red,128; green,128; blue,0}, line width=2pt, short] (13.75,1.5) -- (15,1.5);
			\draw [ color={rgb,255:red,128; green,128; blue,0}, line width=2pt, short] (-6.25,5.25) -- (27.5,5.25);
			\draw [ color={rgb,255:red,128; green,128; blue,0}, line width=2pt, short] (27.5,5.25) -- (27.5,1.5);
			\draw [ color={rgb,255:red,128; green,128; blue,0}, line width=2pt, short] (27.5,1.5) -- (30,1.5);
			\draw [ color={rgb,255:red,255; green,128; blue,0}, line width=2pt, short] (-5,10.25) -- (-5,12.75);
			\draw [ color={rgb,255:red,255; green,128; blue,0}, line width=2pt, short] (-5,12.75) -- (-6.25,12.75);
			\draw [ color={rgb,255:red,255; green,128; blue,0}, line width=2pt, short] (10,11.5) -- (-6.25,11.5);
			\draw [ color={rgb,255:red,255; green,128; blue,0}, line width=2pt, short] (-6.25,10.25) -- (25,10.25);
			\node [font=\Huge, color={rgb,255:red,255; green,128; blue,0}] at (-11.75,12.75) {\textbf{5V/Relais.1}};
			\node [font=\Huge, color={rgb,255:red,255; green,128; blue,0}] at (-11.75,11.5) {\textbf{5V/Relais.2}};
			\node [font=\Huge, color={rgb,255:red,255; green,128; blue,0}] at (-11.75,10.25) {\textbf{5V/Relais.3
			}};
			\node [font=\Huge, color={rgb,255:red,128; green,128; blue,0}] at (-10.5,7.75) {\textbf{D2}};
			\node [font=\Huge, color={rgb,255:red,128; green,128; blue,0}] at (-10.5,6.5) {\textbf{D3}};
			\node [font=\Huge, color={rgb,255:red,128; green,128; blue,0}] at (-10.5,5.25) {\textbf{D4}};
			\node [font=\Huge] at (-13,-22.25) {\textbf{GND(5V)/Relais.1}};
			\node [font=\Huge] at (-13,-23.5) {\textbf{GND(5V)/Relais.2}};
			\node [font=\Huge] at (-13,-24.75) {\textbf{GND(5V)/Relais.3}};
			\draw [ line width=2pt](2.5,3.25) to[short, -o] (2.5,4.25) ;
			\draw [ line width=2pt](5,3.25) to[short, -o] (5,4.25) ;
			\draw [ line width=2pt](7.5,3.25) to[short, -o] (7.5,4.25) ;
			\node [font=\Huge] at (2.5,2.75) {\textbf{NO}};
			\node [font=\Huge] at (7.5,2.75) {\textbf{NC}};
			\node [font=\Huge] at (5,2.75) {\textbf{COM}};
			\draw [ line width=2pt](17.5,3.25) to[short, -o] (17.5,4.25) ;
			\draw [ line width=2pt](20,3.25) to[short, -o] (20,4.25) ;
			\draw [ line width=2pt](22.5,3.25) to[short, -o] (22.5,4.25) ;
			\node [font=\Huge] at (17.5,2.75) {\textbf{NO}};
			\node [font=\Huge] at (22.5,2.75) {\textbf{NC}};
			\node [font=\Huge] at (20,2.75) {\textbf{COM}};
			\draw [ line width=2pt](32.5,3.25) to[short, -o] (32.5,4.25) ;
			\draw [ line width=2pt](35,3.25) to[short, -o] (35,4.25) ;
			\draw [ line width=2pt](37.5,3.25) to[short, -o] (37.5,4.25) ;
			\node [font=\Huge] at (32.5,2.75) {\textbf{NO}};
			\node [font=\Huge] at (37.5,2.75) {\textbf{NC}};
			\node [font=\Huge] at (35,2.75) {\textbf{COM}};
			\draw [ color={rgb,255:red,128; green,0; blue,255}, , line width=2pt](-8.75,17.75) to[short] (35,17.75);
			\draw [ color={rgb,255:red,128; green,0; blue,255}, , line width=2pt](35,17.75) to[short] (35,4.5);
			\draw [ color={rgb,255:red,128; green,0; blue,255}, , line width=2pt](20,4.5) to[short] (20,17.75);
			\draw [ color={rgb,255:red,128; green,0; blue,255}, , line width=2pt](5,4.5) to[short] (5,17.75);
			\node at (5,17.75) [circ, color={rgb,255:red,128; green,0; blue,255}] {};
			\node at (20,17.75) [circ, color={rgb,255:red,128; green,0; blue,255}] {};
			\draw [ color={rgb,255:red,255; green,0; blue,0}, , line width=2pt](32.5,4.5) to[short] (32.5,6.5);
			\draw [ color={rgb,255:red,255; green,0; blue,0}, , line width=2pt](32.5,6.5) to[short] (28.75,6.5);
			\draw [ color={rgb,255:red,255; green,0; blue,0}, , line width=2pt](28.75,6.25) to[short] (28.75,-8.5);
			\draw [ color={rgb,255:red,255; green,0; blue,0}, , line width=2pt](28.75,-8.5) to[short] (33.75,-8.5);
			\draw [ color={rgb,255:red,255; green,0; blue,0}, , line width=2pt](33.75,-8.5) to[short] (33.75,-9.75);
			\draw [ color={rgb,255:red,255; green,0; blue,0}, , line width=2pt](28.75,6) to[short] (28.75,6.5);
			\draw [ color={rgb,255:red,255; green,0; blue,0}, , line width=2pt](17.5,4.25) to[short] (17.5,7.75);
			\draw [ color={rgb,255:red,255; green,0; blue,0}, , line width=2pt](17.5,7.75) to[short] (12.5,7.75);
			\draw [ color={rgb,255:red,255; green,0; blue,0}, , line width=2pt](12.5,7.75) to[short] (12.5,-8.5);
			\draw [ color={rgb,255:red,255; green,0; blue,0}, , line width=2pt](12.5,-8.5) to[short] (18.75,-8.5);
			\draw [ color={rgb,255:red,255; green,0; blue,0}, , line width=2pt](18.75,-8.5) to[short] (18.75,-9.75);
			\draw [ color={rgb,255:red,255; green,0; blue,0}, , line width=2pt](2.5,4.25) to[short] (2.5,7.75);
			\draw [ color={rgb,255:red,255; green,0; blue,0}, , line width=2pt](2.5,7.75) to[short] (-1.25,7.75);
			\draw [ color={rgb,255:red,255; green,0; blue,0}, , line width=2pt](-1.25,7.75) to[short] (-1.25,-8.5);
			\draw [ color={rgb,255:red,255; green,0; blue,0}, , line width=2pt](-1.25,-8.5) to[short] (3.75,-8.5);
			\draw [ color={rgb,255:red,255; green,0; blue,0}, , line width=2pt](3.75,-8.5) to[short] (3.75,-9.75);
			\draw [ line width=2pt](3.75,-11) to[short, -o] (3.75,-10) ;
			\draw [ line width=2pt](6.25,-11) to[short, -o] (6.25,-10) ;
			\draw [ line width=2pt](18.75,-11.25) to[short, -o] (18.75,-10) ;
			\draw [ line width=2pt](21.25,-11.25) to[short, -o] (21.25,-10) ;
			\draw [ line width=2pt](33.75,-11.25) to[short, -o] (33.75,-10) ;
			\draw [ line width=2pt](36.25,-11.25) to[short, -o] (36.25,-10) ;
			\node [font=\Huge, color={rgb,255:red,128; green,0; blue,255}] at (-11.25,17.75) {\textbf{+ 12 V}};
			\draw [ line width=2pt](-8.75,-29.75) to[short] (40,-29.75);
			\draw [ line width=2pt](40,-29.75) to[short] (40,-8.5);
			\draw [ line width=2pt](40,-8.5) to[short] (36.25,-8.5);
			\draw [ line width=2pt](36.25,-8.5) to[short] (36.25,-10);
			\draw [ line width=2pt](26.25,-24.25) to[short] (26.25,-24.75);
			\draw [ line width=2pt](21.25,-10) to[short] (21.25,-8.5);
			\draw [ line width=2pt](21.25,-8.5) to[short] (25,-8.5);
			\draw [ line width=2pt](25,-8.5) to[short] (25,-29.75);
			\draw [ line width=2pt](6.25,-9.75) to[short] (6.25,-8.5);
			\draw [ line width=2pt](6.25,-8.5) to[short] (10,-8.5);
			\draw [ line width=2pt](10,-8.5) to[short] (10,-29.75);
			\node at (10,-29.75) [circ] {};
			\node at (25,-29.75) [circ] {};
			\node [font=\Huge] at (-12,-29.75) {\textbf{GND (12V)}};
			\node [font=\Huge] at (4.75,-12) {\textbf{+    -}};
			\node [font=\Huge] at (19.75,-12.25) {\textbf{+    -}};
			\node [font=\Huge] at (35,-12) {\textbf{+    -}};
		\end{circuitikz}
	}%
		\label{12VSchaltplan}
\end{figure}

\section{Software-Schnittstellen}
In dem Projekt Magic Faucet Fountain, wird die nRF Connect-App als Software-Schnittstelle zum Arduino verwendet. Diese kommuniziert über BLE mit dem Arduino. Mit der App kann man den Magic Faucet Fountain steuern, Messdaten empfangen und Statusmeldungen abgreifen.

\subsubsection{Bluetooth-Kommunikationsschnittstelle}
Die an der Bluetooth-Kommunikation beteiligten Komponenten sind der Arduino mit seinem integrierten BLE-Modul und das Smartphone mit der nRF-Connect-App. 

\subsubsection{Arduino-Programmierschnittstelle}
Der Arduino Mikrokontroller wird mit einem Arduino Sketch (Arduino IDE) programmiert, der folgende Aufgaben übernimmt:
\begin{itemize}
	\item Initialisierung der Pins
	\item Einrichtung der seriellen Schnittstelle (Serial.begin(9600))
	\item Einlesen und Interpretieren von Bluetooth-Befehlen
	\item Setzen der Ausgänge zur Ansteuerung der Aktoren
\end{itemize}

\subsubsection{Anwendungsprogrammierschnittstelle auf Befehlsebene}
Die Bluetooth-App kommuniziert mit dem Arduino über eine einfache Befehlssprache, z.B.:
\begin{itemize}
	\item „EIN“: Pumpe ein
	\item „AUS“: Pumpe aus
\end{itemize}
Der Arduino-Sketch interpretiert diese Befehle als Teil eines einfachen protokollbasierten Interfaces. Diese Art der Schnittstelle ist textbasiert.

\subsection{Ablaufdiagramm der Anwendung}
Zur Veranschaulichung des Zusammenspiels der beschriebenen Software-Schnittstellen dient das in Abbildung \ref{ablaufdiagramm} dargestellte Ablaufdiagramm. Es zeigt den strukturierten Programmablauf vom Start der Anwendung über die Initialisierung und Bluetooth-Kommunikation bis hin zur Ansteuerung der LEDs und der Wasserpumpe. Zudem wird der Umgang mit kritischen Zuständen wie zu hohen oder zu niedrigen Wasserständen abgebildet. 
\begin{figure}[H]
	\centering
	\caption{Ablaufdiagramm}
	\label{ablaufdiagramm}
	\resizebox{1\textwidth}{!}{%
		\begin{circuitikz}
			\tikzstyle{every node}=[font=\Huge]
			\draw [ line width=2pt , rounded corners = 12.0] (-53.25,175.25) rectangle  node {\Huge \textbf{Start nRF Connect}} (-43.75,172.5);
			\draw [ line width=2pt , rounded corners = 12.0] (-58.75,180.25) rectangle  node {\Huge \textbf{Start}} (-51.25,177.75);
			\draw [ line width=2pt , rounded corners = 12.0] (-65,175.25) rectangle  node {\Huge \textbf{Programm Initialisieren}} (-55,172.75);
			\draw [ color={rgb,255:red,0; green,0; blue,255} , line width=2pt , rounded corners = 12.0] (-63.75,170.25) rectangle  node {\Huge \textbf{Grüne LED EIN}} (-56.25,167.75);
			\draw [line width=2pt, short] (-51.25,179) -- (-47.5,179);
			\draw [line width=2pt, short] (-58.75,179) -- (-60,179);
			\draw [line width=2pt, short] (-60,179) -- (-60,177.75);
			\draw [line width=2pt, ->, >=Stealth] (-60,177.75) -- (-60,175.5);
			\draw [line width=2pt, ->, >=Stealth] (-47.5,179) -- (-47.5,175.5);
			\draw [line width=2pt, ->, >=Stealth] (-60,172.75) -- (-60,170.5);
			\draw [ line width=2pt , rounded corners = 12.0] (-51,166.375) -- (-47.375,164.5) -- (-43.75,166.375) -- (-47.375,168.25) -- cycle;
			\node [font=\Huge] at (-47.375,166.375) {\textbf{BT Signal?}};
			\draw [line width=2pt, short] (-43.75,166.5) -- (-40.5,166.5);
			\draw [ color={rgb,255:red,0; green,0; blue,255} , line width=2pt , rounded corners = 12.0] (-40,167.75) rectangle  node {\Huge \textbf{Blaue LED Blinken}} (-31.25,165.25);
			\draw [line width=2pt, short] (-36.5,167.75) -- (-36.5,169.5);
			\draw [line width=2pt, short] (-36.5,169.5) -- (-36.5,170.25);
			\draw [line width=2pt, short] (-47.5,172.5) -- (-47.5,168.25);
			\draw [line width=2pt, ->, >=Stealth] (-47.5,169.5) -- (-47.5,168.25);
			\draw [line width=2pt, ->, >=Stealth] (-36.5,170.25) -- (-47.25,170.25);
			\draw [ color={rgb,255:red,0; green,0; blue,255} , line width=2pt , rounded corners = 12.0] (-51.25,161.5) rectangle  node {\Huge \textbf{Blaue LED EIN}} (-43.75,159);
			\draw [line width=2pt, ->, >=Stealth] (-47.5,164.5) -- (-47.5,161.5);
			\draw [line width=2pt, ->, >=Stealth] (-41.25,166.5) -- (-40.25,166.5);
			\node [font=\Huge] at (-42.25,167.25) {\textbf{Nein}};
			\node [font=\Huge] at (-48.5,163.25) {\textbf{Ja}};
			\draw [ line width=2pt , rounded corners = 12.0] (-67,161.875) -- (-62.375,159.75) -- (-57.75,161.875) -- (-62.375,164) -- cycle;
			\node [font=\Huge] at (-62.375,161.875) {\textbf{Pegel < 20 cm?}};
			\draw [line width=2pt, short] (-60,167.75) -- (-60,166.5);
			\draw [line width=2pt, short] (-60,166.5) -- (-62.25,166.5);
			\draw [line width=2pt, ->, >=Stealth] (-62.25,166.5) -- (-62.25,164);
			\draw [line width=2pt, ->, >=Stealth] (-62.5,159.75) -- (-62.5,156.75);
			\node [font=\Huge] at (-63.25,158.25) {\textbf{Ja}};
			\draw [ color={rgb,255:red,0; green,0; blue,255} , line width=2pt , rounded corners = 12.0] (-66.25,156.5) rectangle  node {\Huge \textbf{Pumpe EIN}} (-58.75,154);
			\draw [line width=2pt, ->, >=Stealth] (-62.5,154) -- (-62.5,152.75);
			\draw [ line width=2pt , rounded corners = 12.0] (-67,151) -- (-62.625,149.25) -- (-58.25,151) -- (-62.625,152.75) -- cycle;
			\node [font=\Huge] at (-62.625,151) {\textbf{Pegel > 20 cm?}};
			\draw [line width=2pt, ->, >=Stealth] (-66.75,151) -- (-70,151);
			\draw [ color={rgb,255:red,255; green,0; blue,0} , line width=2pt , rounded corners = 12.0] (-84.75,150.25) rectangle  node {\Huge \textbf{Warnung: Pegel Maximum erreicht!}} (-70,151.5);
			\node [font=\Huge] at (-68.25,151.75) {\textbf{Ja}};
			\draw [ color={rgb,255:red,0; green,0; blue,255} , line width=2pt , rounded corners = 12.0] (-81.75,147.75) rectangle  node {\Huge \textbf{Pumpe AUS}} (-73.75,145.25);
			\draw [line width=2pt, ->, >=Stealth] (-77.5,150.25) -- (-77.5,148);
			\draw [line width=2pt, short] (-58.25,151) -- (-56.25,151);
			\draw [line width=2pt, short] (-56.25,151) -- (-56.25,155.25);
			\draw [line width=2pt, ->, >=Stealth] (-56.25,155.25) -- (-58.5,155.25);
			\node [font=\Huge] at (-57.5,151.75) {\textbf{Nein}};
			\draw [line width=2pt, short] (-57.75,161.75) -- (-54.75,161.75);
			\draw [line width=2pt, short] (-55,161.75) -- (-55,146.5);
			\draw [line width=2pt, ->, >=Stealth] (-55,146.5) -- (-73.5,146.5);
			\node [font=\Huge] at (-56.5,162.5) {\textbf{Nein}};
			\draw [ line width=2pt , rounded corners = 12.0] (-78.75,162) -- (-74.25,160) -- (-69.75,162) -- (-74.25,164) -- cycle;
			\node [font=\Huge] at (-74.25,162) {\textbf{Pegel < 5 cm?}};
			\draw [line width=2pt, short] (-74.25,160) -- (-74.25,155.5);
			\draw [line width=2pt, ->, >=Stealth] (-74.25,155.25) -- (-66.5,155.25);
			\node [font=\Huge] at (-75.5,158) {\textbf{Nein}};
			\draw [line width=2pt, ->, >=Stealth] (-78.75,162) -- (-80.25,162);
			\draw [ color={rgb,255:red,255; green,0; blue,0} , line width=2pt , rounded corners = 12.0] (-95.5,162.75) rectangle  node {\Huge \textbf{Warnung: Pegel Minimum erreicht!}} (-80.25,161.25);
			\node [font=\Huge] at (-78.5,163.25) {\textbf{Ja}};
			\draw [line width=2pt, short] (-83.75,161) -- (-83.75,155.25);
			\draw [line width=2pt, ->, >=Stealth] (-83.75,155.25) -- (-74.5,155.25);
			\draw [line width=2pt, ->, >=Stealth] (-66.75,162) -- (-69.75,162);
		\end{circuitikz}
	}%
\end{figure}

\subsection{Code der Anwendung}

\begin{code}[H]
	\label{code:codeApplikation}
	\ArduinoExternal{}{../../Code/Arduino/Applikation/ablaufApplikation/ablaufApplikation.ino}
\end{code}

\section{Installation}
Die Installation des Magic Faucet Fountain gliedert sich in zwei Hauptbereiche: den mechanischen Aufbau der Baugruppe und die elektrische Inbetriebnahme der Steuerungskomponenten. Im Folgenden werden beide Teile Schritt für Schritt erläutert.

\subsection{Mechanischer Aufbau}
\begin{itemize}
	\item Der Blumentopf wird als zentrales Bauteil aufgestellt und mit einem speziell angefertigten Deckel verschlossen. Dieser beinhaltet:
	\begin{itemize}
        	\item Eine zentrale Öffnung zur Durchführung des Acrylrohrs, das den Wasserfluss führt.
        	\item Mehrere Rücklaufbohrungen, durch die Wasser in den Topf zurückfließt.
	        \item Eine Gewindeaufnahme auf der Unterseite zur Befestigung des Führungsrohrs für den Schwimmer.
	\end{itemize}
	\item Das transparente Acrylrohr wird vertikal durch den Deckel geführt und oben mit einem dekorativen Wasserhahn verbunden. Kleine Bohrungen im Rohr direkt unter dem Hahn ermöglichen den sichtbaren Wasseraustritt.
	\item Das Führungsrohr wird senkrecht unter dem Deckel montiert. Im Inneren befindet sich ein Schwimmer, der sich mit dem Wasserstand bewegt.
	\item Der untere Abschluss des Führungsrohrs besteht aus einem fest verschlossenen Deckel mit integriertem Gitter, das den Schwimmer sichert und den Wassereintritt erlaubt.
	\item Der Ultraschallsensor wird oberhalb des Führungsrohrs an der Aufnahme befestigt und misst den Abstand zum Schwimmer zur Bestimmung des Füllstands.
	\item Auf dem Deckel können zur Gestaltung Dekosteine oder andere Elemente platziert werden.
\end{itemize}

\subsection{Elektrische Installation}
Alle elektronischen Komponenten sind im separaten E-Kasten untergebracht, der außerhalb des Blumentopfs montiert wird und Feuchtigkeitsschutz sowie einfache Wartung ermöglicht.

Folgende Komponenten sind dort integriert:
\begin{itemize}
    \item Arduino-Mikrocontroller, als zentrale Steuereinheit.
    \item Relais-Schaltkreis zur Ansteuerung der Tauchpumpe, die das Wasser durch das Acrylrohr fördert.
    \item Ultraschallsensor, über digitale Pins mit dem Arduino verbunden.
    \item Zwei Status-LEDs:
    	\begin{itemize}
        \item Grüne LED: zeigt den Betriebszustand an – leuchtet, wenn der Kippschalter eingeschaltet ist.
        \item Blaue LED: zeigt den Bluetooth-Zustand – blinkt bei Verbindungssuche, leuchtet bei bestehender Verbindung.
        \end{itemize}
    \item Kippschalter als Hauptschalter für das gesamte System.
    \item Stromversorgung über externes Netzteil, angeschlossen an den E-Kasten.
    \item Alle Verbindungen zur Pumpe, zu den Sensoren und zu den LEDs sind über robuste Steckverbindungen aus dem E-Kasten herausgeführt.
\end{itemize}

\subsection{Software}
Damit der Magic Faucet Fountain vollständig genutzt werden kann, ist die App \textit{nRF Connect} zwingend erforderlich. Stellen Sie sicher, dass diese auf Ihrem Smartphone installiert ist, mehr Infos im Kapitel \ref{chap:einstiegNRF}. Sobald diese installiert ist, kann die Konfiguration und Inbetriebnahme des Magic Faucet Fountain beginnen.

\section{Konfiguration und Inbetriebnahme}
Die Konfiguration des Magic Faucet Fountain umfasst die Inbetriebnahme und Abstimmung der elektronischen Komponenten sowie die Kopplung mit der mobilen App. Ziel ist es, die Steuerung des Brunnens per Bluetooth zuverlässig zu ermöglichen und eine präzise Füllstandsanzeige über den Ultraschallsensor bereitzustellen.

\subsection{Elektrische Inbetriebnahme}
\begin{itemize}
	\item Vor dem Einschalten müssen alle elektrischen Verbindungen überprüft werden: Die Pumpe, der Ultraschallsensor, die Status-LEDs, die Relais, der Kippschalter und das Bluetooth-Modul müssen korrekt mit dem Arduino verbunden sein.
	\item Die Stromversorgung wird über den Kippschalter aktiviert. Beim Einschalten leuchtet die grüne Betriebs-LED zur Bestätigung.
\end{itemize}

\subsection{Installation und Konfiguration mit nRF Connect}
\textbf{Hier fehlen noch Screenshots wie man die App einrichtet und benutzt mit der Applikation}

\subsection*{Verbindung mit der Applikation herstellen}

\subsection*{App-Ansicht}

\subsection*{Steuerung des Magic Faucet Fountain}

\begin{itemize}
	\item Nach dem Einschalten beginnt die blaue LED zu blinken, was signalisiert, dass Bluetooth aktiv ist und auf eine Verbindung wartet.
	\item Über die App wird das Bluetooth-Modul des Brunnens gesucht und gekoppelt. Nach erfolgreicher Verbindung wechselt die blaue LED von Blinken auf Dauerlicht.
	\item Die App ermöglicht das manuelle Ein- und Ausschalten der Pumpe sowie die Anzeige des aktuellen Wasserstands.
\end{itemize}



\section{Einschränkungen}
Da das Projekt sich noch in der finalen Planungs- und Dokumentationsphase befindet, konnten einige Einschränkungen nur theoretisch erkannt werden:
\begin{itemize}
	\item Praxistests fehlen bislang, weshalb reale Probleme bei der Wasserführung, Abdichtung oder elektrischen Sicherheit noch nicht abschließend beurteilt werden können.
	\item Die Stromversorgung muss bei der Kombination von 5 V- und 12 V-Komponenten mit Bedacht gewählt werden, um Störungen oder Überlastungen zu vermeiden.
	\item Die Bluetooth-Kommunikation hängt von der Stabilität der Verbindung und der verwendeten App ab, mögliche Verbindungsabbrüche sind noch nicht erprobt.
	\item Platzbedarf und Montage könnten im späteren Aufbau auf bisher nicht eingeplante bauliche Herausforderungen stoßen.
\end{itemize}



